\documentclass[conference]{IEEEtran}
\usepackage[spanish]{babel}
\usepackage[utf8]{inputenc}
\usepackage[T1]{fontenc}
\usepackage{booktabs}
\usepackage{url}

\title{Parametros, Marcos de Activacion y Pila de Ejecucion}
\author{\IEEEauthorblockN{Ronquillo Nunez Braulio}
\IEEEauthorblockA{Analisis y Diseno de Algoritmos}}

\begin{document}
\maketitle

\begin{abstract}
El paso de parametros en funciones puede analizarse mediante una pila de
ejecucion. Cada llamada crea un marco de activacion que contiene parametros,
variables locales y control de retorno. Este reporte describe la estructura
LIFO, su relacion con llamadas anidadas y recursivas, y la simulacion
implementada para observar operaciones de llamada, asignacion y retorno.
\end{abstract}

\begin{IEEEkeywords}
parametros, pila, LIFO, marco de activacion, recursion.
\end{IEEEkeywords}

\section{Introduccion}
Cuando un programa invoca una funcion, el flujo de ejecucion se suspende en
el punto de llamada y se transfiere control a otro bloque de codigo. Para
recuperar el estado anterior se utiliza una estructura conceptual conocida
como pila de ejecucion. Esta pila almacena informacion de cada llamada activa.

NIST define una pila como una coleccion donde solo puede retirarse el elemento
agregado mas recientemente; sus operaciones basicas son \texttt{push} y
\texttt{pop}, y se conoce como LIFO \cite{nist_stack}. Esta propiedad encaja
con las llamadas de funcion: la ultima funcion invocada es la primera que debe
terminar.

\section{Marco Teorico}
Un marco de activacion representa una llamada particular. Aunque los detalles
dependen del lenguaje y compilador, el modelo abstracto incluye:
\begin{itemize}
\item parametros de entrada,
\item variables locales,
\item direccion de retorno,
\item informacion temporal para evaluar expresiones.
\end{itemize}
En llamadas anidadas, cada nuevo marco se apila sobre el anterior. Cuando la
funcion termina, se desapila su marco y el control regresa al llamador.

Este mecanismo tambien explica la recursion. Cada llamada recursiva posee sus
propias variables locales, aunque ejecute el mismo codigo. El limite practico
esta dado por el tamano maximo de pila; si se excede, aparece un error de
desbordamiento.

\section{Modelo del Programa}
La simulacion implementada usa tres operaciones principales:
\begin{table}[htbp]
\caption{Operaciones simuladas}
\centering
\begin{tabular}{ll}
\toprule
Operacion & Efecto \\
\midrule
\texttt{CALL f} & inserta un marco de funcion \\
\texttt{SET x v} & asigna una variable local \\
\texttt{RETURN} & elimina el marco superior \\
\bottomrule
\end{tabular}
\end{table}

La pila se representa con una lista. El extremo final de la lista funciona
como tope. Esta decision permite que insertar y retirar el marco superior sea
una operacion directa. El programa reporta la evolucion de la pila y detecta
acciones invalidas, como intentar retornar cuando no existe una llamada activa.

\section{Criterios de Calidad}
El ejercicio permite revisar tres criterios de calidad:
\begin{enumerate}
\item \textit{Orden correcto}: la estructura debe conservar el orden LIFO.
\item \textit{Aislamiento}: una variable local solo debe afectar el marco
actual.
\item \textit{Control de tamano}: la simulacion debe advertir si se supera la
profundidad maxima permitida.
\end{enumerate}
Estos criterios son relevantes porque un error en cualquiera de ellos cambia
el significado de las llamadas de funcion.

\section{Complejidad}
Para $q$ operaciones, el costo total es $O(q)$ si se omite el costo de imprimir
la pila completa. El espacio usado es $O(h)$, donde $h$ es la maxima
profundidad de llamadas activas. En el peor caso, si nunca se ejecuta
\texttt{RETURN}, $h=q$.

\section{Conclusion}
La pila de ejecucion ofrece una forma clara de estudiar el paso de parametros.
El comportamiento LIFO garantiza que cada funcion termine antes de regresar al
contexto que la invoco. La simulacion ayuda a visualizar recursividad,
variables locales y errores de control de llamadas.

\begin{thebibliography}{00}
\bibitem{nist_stack} P. E. Black, ``stack,'' Dictionary of Algorithms and Data Structures, NIST, 2019. [Online]. Available: \url{https://www.nist.gov/dads/HTML/stack.html}
\bibitem{mit_algorithms} E. Demaine, J. Ku, and J. Solomon, ``6.006 Introduction to Algorithms: Lecture Notes,'' MIT OpenCourseWare, 2020. [Online]. Available: \url{https://ocw.mit.edu/courses/6-006-introduction-to-algorithms-spring-2020/pages/lecture-notes/}
\bibitem{ods_stack} P. Morin, ``Open Data Structures,'' Athabasca University Press, 2013. [Online]. Available: \url{https://opendatastructures.org/}
\end{thebibliography}

\end{document}
